% Round 15. New assembly, requiring separate independent review.
% This file does not modify any Round 1--14 result.
\section{PWS integration along a projective line over an arbitrary smooth base}
\label{r15hp:section}

The local perturbation of presenting moments changes the scope of the
relative construction.  It permits a new splitting tuple at each base
point, rather than requiring one fixed global tuple to remain separable.
The result below is a pushforward for the particular projection with
fiber $\mathbf P^1$, on PWS inputs.  It is not a proper pushforward
theorem for arbitrary actual delta-forms.

We work over the original finite valued extension $K/\mathbf Q_p$.
Write $\operatorname{PS}^{p,q}=\mathcal B^{p,q,0}$, with the
ordinary topology throughout.  An ordinary open in a smooth algebraic
analytification is boundaryless and has the same pure dimension as its
ambient component.  The notation $U\times\mathbf P^1$ means the
analytic inverse image of $U$ under the algebraic projection; no
Cartesian-product topology is asserted.

\begin{theorem}[Ordinary PWS integration along $\mathbf P^1$]
\label{r15hp:projection}
Let $S/K$ be a smooth integral algebraic variety of dimension $d$,
let $U\subset S^{\mathrm{an}}$ be an ordinary open, and let
$\pi:U\times\mathbf P^1\to U$ be the projection.
For $p,q\geq1$ and every
$\omega\in\operatorname{PS}^{p,q}(U\times\mathbf P^1)$,
there is a unique ordinary PWS form
\[
 \mathcal I_U\omega\in\operatorname{PS}^{p-1,q-1}(U)
\tag{HP15.1}\label{r15hp:output}
\]
such that
\[
 \int_U\mathcal I_U\omega\wedge\eta
       =\int_{U\times\mathbf P^1}\omega\wedge\pi^*\eta
\tag{HP15.2}\label{r15hp:all-smooth-pairing}
\]
for every ordinary smooth compactly supported test
$\eta\in\mathcal A_c^{d+1-p,d+1-q}(U)$.
The operator is linear and commutes with restriction to ordinary
base opens.  Degrees outside the dimension range are interpreted
as zero; if $p=0$ or $q=0$, the current pushforward is zero by
relative degree.  Every auxiliary finite cover is removed using
its full native algebraic degree.
\end{theorem}

\begin{proof}
We construct the representative near an arbitrary $x\in U$.

\emph{Complete-fiber presentations and a finite etale splitting.}
Choose a smooth affine algebraic chart $V$ containing $x$.
Lemma~\ref{r15b:Laurent-fiber-exists} applies to $\omega$ on the
inverse image of $U\cap V^{\mathrm{an}}$.  It supplies finitely
many algebraic Laurent source domains whose ordinary interiors
cover the complete fiber over $x$, and then cover the complete
inverse image of a smaller ordinary base neighborhood.  Their
defining conditions have scalar constant thresholds on compact
coordinate boxes.  All presenting units are rationalized as
quotients of polynomials in $z$ or $1/z$ with coefficients regular
on $V$.

Apply Proposition~\ref{r15b:separable-perturbation}.  Its changes
preserve each entire source domain, including its boundary and
every point excluded from it, and preserve every presenting unit
norm on an ordinary neighborhood of the domain.  Consequently
the represented actual form is unchanged.  At $x$, the new
polynomials have nonzero leading and constant coefficients and
pairwise disjoint simple roots.  Corollary~\ref{r15b:local-etale-splitting}
therefore gives a Zariski open $V^g$ containing $x$ and a finite
etale surjective map
\[
                    a:V'\longrightarrow V^g
\tag{HP15.3}\label{r15hp:local-cover}
\]
of smooth integral algebraic varieties on which all those roots
split as regular units.  Their pairwise differences are units.
Append the fixed marks $0$ and $\infty$.

If needed, the cover may be chosen Galois: order all roots of the
total monic polynomial in the finite etale splitting algebra and
choose a connected component.  Its function field is generated
by all the roots, hence is the splitting field of the original
polynomial.  It is a finite Galois extension.  The resulting
normal finite etale cover is a torsor for that Galois group.
No ramified normalization or alteration is used here.

Choose a compact affinoid neighborhood $B$ of $x$ inside the
ordinary neighborhood where the complete-fiber cover holds, and
put $B'=a^{-1}(B)$.  Set $\Omega=B^\circ$ and
$\Omega'=a^{-1}(\Omega)$, shrinking $B$ first if necessary.
The compact affinoid $B'$ contains $\Omega'$ in its ordinary
interior.  The pulled-back form is denoted $\omega'$.  The
finite base tuple $b$ contains all leading factors and all
nonzero root differences.  Proposition~\ref{r15b:local-whole-source-presentation}
gives one intrinsic PWS tensor for $\omega'$ on the whole finite
polyhedral moment image of
\[
 W_{B'}=(B'\times\mathbf A^1_z)\setminus
                         \bigcup_{i=1}^m\{z=a_i\},
 \qquad a_1=0,
\tag{HP15.4}\label{r15hp:whole-source}
\]
in the genuine tuple $(b,z-a_1,\ldots,z-a_m)$.  All source-domain
inequalities, including the fixed coordinate boxes, are accounted
for in this tuple.  This is a statement on the complete moment
image over $B'$, not only on a skeleton or on selected fiber
points.

\emph{Excision of every marked tail.}
For a finite marked section $s_i$ or for the infinity section,
the actual form
\[
              \omega'-\pi'^*s_i^*\omega'
\tag{HP15.5}\label{r15hp:graph-difference}
\]
has zero restriction to that algebraic graph.  The frozen
algebraic immersion-germ isomorphism, Round 14
\texttt{r14l:germ-iso}, makes it zero on an ordinary neighborhood
of the graph.  It applies locally even though $\omega'$ is only
given on an ordinary open.  The graph over $B'$ is compact.
The continuous chordal distance to it has a strictly positive
minimum on the complement of any chosen such neighborhood in
$B'\times\mathbf P^1$.  Thus one can choose a uniform small
chordal tube on which
$\omega'=\pi'^*\beta_i$, where $\beta_i=s_i^*\omega'$.

Choose a scalar PWS cutoff $\theta_i$ which is zero on a smaller
tube and one off that tube.  On the complement of the section it
can be taken as a smooth function of
\[
 v(z-a_i)-\min(0,v(z))-\min(0,v(a_i)),
\tag{HP15.6}\label{r15hp:chordal-moments}
\]
and for infinity as a smooth function of $-\min(0,v(z))$.
For the section $0$ the term $\min(0,v(a_i))$ is simply zero.
The cutoff is constant zero near the deleted graph, so extends
there as an ordinary PWS scalar.  All its moments belong to the
same tuple: $z=z-a_1$, and $a_i=a_i-a_1$ is among the nonzero
base differences when $i\ne1$.  Put $\theta=\prod_i\theta_i$,
including the infinity factor.

For any smooth compact complementary test $\eta'$ on $\Omega'$,
the discarded term on the $i$th tube satisfies
\[
 \omega'\wedge\pi'^*\eta'
     =\pi'^*(\beta_i\wedge\eta')=0.
\tag{HP15.7}\label{r15hp:tail-degree}
\]
Indeed the form inside the pullback has bidegree $(d+1,d+1)$
on a $d$-dimensional space.  This is equality of actual forms,
before integration.  Applying the elementary product-cutoff
identity to $1-\theta$ now proves
\[
 \pi'_*[\theta\omega']=\pi'_*[\omega']
\tag{HP15.8}\label{r15hp:tail-current}
\]
against every prescribed smooth test.  No derivative of a cutoff
is introduced, and no vertical derivative or fiberwise value is
being substituted for this degree calculation.

\emph{A genuine finite presentation with compact moment support.}
The retained form $\theta\omega'$ is zero on ordinary
neighborhoods of every deleted graph and infinity.  Its support
in the complement of those graphs is proper over $\Omega'$.
The uniform lower chordal bounds and the infinity cutoff bound
all $|z-a_i|$ above and away from zero over $B'$.  The base
moments are bounded on $B'$.  Hence its entire moment support is
contained in a compact portion of the finite polyhedral image
in \eqref{r15hp:whole-source}.

Multiplication of that intrinsic tensor by the displayed cutoffs
gives one PWS tensor on the entire source image; it is tangentially
zero off that compact portion.  For precision, take three nested
large Euclidean cubes, with the support in the interior of the
smallest.  Extend the intrinsic tensor near its intersection with
the largest cube using the frozen degreewise finite-polyhedron
PWS extension lemma, Round 14 \texttt{r14f:extension}.  Multiply
the extension by a Euclidean smooth cutoff equal to one near the
smallest intersection and supported in the extension neighborhood
and in the interior of the largest cube.  It can be extended by
zero to the entire ambient Euclidean space.  On the source image
it agrees with $\theta\omega'$ everywhere: in the intermediate
collar the original tangential tensor is zero.  Thus this is one
finite ordinary PWS presentation in $(b,z-a_i)$ on
$W_{\Omega'}$.  It is not zero extension of a tensor restricted
by the indicator of a moving polyhedral band.

\emph{All mixed tests on the splitting cover.}
Proposition~\ref{r15ray:analytic} now applies, with base
$\Omega'$ of dimension $d$ and input degree $(p,q)$.
The uniformly removed tails supply its finite integration bounds.
It produces an ordinary PWS form $\tau'$ of degree
$(p-1,q-1)$ satisfying
\[
 \int_{\Omega'}\tau'\wedge\eta'
      =\int_{\Omega'\times\mathbf P^1}\omega'
                                                \wedge\pi'^*\eta'
\tag{HP15.9}\label{r15hp:upstairs-pairing}
\]
for every ordinary smooth compact test $\eta'$.
The cited proposition adds every moment of that test, proves the
refined skeleton is the marked-tree family over the refined base
skeleton, retains the native base weight times the weight-one
relative edge, and integrates the contracted tensor on every
minimum-label band.  It therefore covers mixed tests and target
moment images of insufficient initial rank.  No scalar-only
fiber calculation is used as a substitute for that assertion.

\emph{Native finite descent and ordinary gluing.}
Let $e=\deg_K a$ be the full native degree, and define on $\Omega$
\[
             \tau=e^{-1}\operatorname{Tr}^{\mathrm{PS}}_a\tau'.
\tag{HP15.10}\label{r15hp:finite-descent}
\]
This uses the ordinary sheaf trace of Round 14
\texttt{r14f:pws-trace} for a finite dominant map of normal
integral algebraic varieties, restricted to $\Omega$.
It is available in every PWS bidegree.  For a prescribed test
$\eta$ on $\Omega$, its all-actual-test pairing and
\eqref{r15hp:upstairs-pairing} give
\[
 \begin{split}
 \int_\Omega\tau\wedge\eta
 &=e^{-1}\int_{\Omega'}\tau'\wedge a^*\eta\\
 &=e^{-1}\int_{\Omega'\times\mathbf P^1}
                  (a\times\mathrm{id})^*
                       (\omega\wedge\pi^*\eta)\\
 &=\int_{\Omega\times\mathbf P^1}\omega\wedge\pi^*\eta.
 \end{split}
\tag{HP15.11}\label{r15hp:native-calibration}
\]
The final equality is the frozen native finite degree formula.
The map on the product has degree $e$, not a higher power of
$e$.  All relevant paired supports are compact by properness of
the projective-line projection.  Constant-field extensions,
if present in the splitting field, are included in $e$.

Local representatives constructed near different points or with
different source domains, perturbations, labels, cutoffs, or
finite covers have the same ordinary current on every ordinary
overlap.  Mihatsch Proposition 4.11 makes them equal as actual
forms.  They glue to \eqref{r15hp:output}; the same injectivity
gives uniqueness, linearity, and restriction compatibility.
\end{proof}

\begin{lemma}[Finite etale Galois compatibility]
\label{r15hp:finite-compatibility}
Let $a:S'\to S$ be a finite etale Galois surjective morphism of
smooth integral algebraic varieties, with group $H$, and restrict
it to an ordinary open $U$ and its inverse image $U'$.
For every PWS input $\omega$ as in Theorem~\ref{r15hp:projection},
\[
 a^*\mathcal I_U\omega
       =\mathcal I_{U'}((a\times\mathrm{id})^*\omega).
\tag{HP15.12}\label{r15hp:finite-pullback}
\]
\end{lemma}

\begin{proof}
Write $\tau'$ for the right side.  In its defining current
pairing, change variables by any $\sigma\in H$.  The input is
$H$-invariant and the automorphism has native degree one.
Current injectivity gives $\sigma^*\tau'=\tau'$ as an actual
PWS form.  The same calculation as
\eqref{r15hp:native-calibration} gives
\[
 \operatorname{Tr}^{\mathrm{PS}}_a\tau'
                              =|H|\mathcal I_U\omega.
\]
The orbit identity in the proof of Round 14
\texttt{r14f:pws-trace}, or equivalently its Galois descent
theorem \texttt{r14f:galois-pws}, gives
\[
 a^*\operatorname{Tr}^{\mathrm{PS}}_a\tau'
             =\sum_{\sigma\in H}\sigma^*\tau'=|H|\tau'.
\]
Divide by the positive native degree $|H|$.
This proof does not assume a trace base-change theorem.
\end{proof}

\begin{proposition}[Algebraic base-change naturality]
\label{r15hp:base-change}
Let $b:T\to S$ be an arbitrary algebraic morphism between smooth
integral $K$-varieties.  It need not be dominant or equidimensional.
Let $U\subset S^{\mathrm{an}}$ and $V\subset T^{\mathrm{an}}$
be ordinary opens with $b(V)\subset U$.  For every PWS input
$\omega$ on $U\times\mathbf P^1$,
\[
 b^*\mathcal I_U\omega
       =\mathcal I_V((b\times\mathrm{id})^*\omega)
\tag{HP15.13}\label{r15hp:base-change-formula}
\]
as ordinary actual PWS forms in the indicated degree.
\end{proposition}

\begin{proof}
The assertion is local on $V$.  Around $x=b(t)$ choose the
finite Galois etale splitting cover and the retained finite
presentation used in the proof of Theorem~\ref{r15hp:projection}.
Pull back the algebraic cover to an algebraic neighborhood of
$t$ in $T$.  The resulting finite etale $H$-torsor may be
disconnected.  Choose a connected component $T'$.  Its stabilizer
$H_0\subset H$ acts transitively on its geometric fibers; thus
$T'\to T$ on that neighborhood is a finite etale Galois cover
with group $H_0$, and is surjective.  It is smooth and integral.
One may equivalently verify this on the finite etale torsor and
its open-and-closed connected components.

On the common splitting cover, pull back the roots, all base
units, and all the cutoffs.  They remain genuine units and
distinct marked sections.  The root-tail equality
\eqref{r15hp:graph-difference} pulls back as an actual equality,
and its excision proof \eqref{r15hp:tail-degree} still applies
with $\dim T$ in place of $d$.  The pulled retained form has the
same finite tensor presentation through the pulled-back moments.

The moment map from any compact target chart on $T'$ to the
original base moment image is a restriction of a linear
coordinate map after adjoining the pulled-back base units.
The distance functions, the root-ray maps, the minimum-label
bounds, and the two vertical contractions pull back by that map.
Lemma~\ref{r15ray:band} therefore identifies the entire band
integral with the pullback of the original band integral,
including on every collapsed or lower-dimensional image face.
Proposition~\ref{r15ray:analytic} identifies both sides with
their prescribed currents against all tests on their respective
bases.  Consequently \eqref{r15hp:base-change-formula} holds on
this common splitting cover as an equality of actual PWS forms.

Lemma~\ref{r15hp:finite-compatibility} identifies the pullbacks
of the two downstairs operators with exactly those two
upstairs operators.  Faithfulness of finite dominant actual
pullback then proves the equality downstairs near $t$.
Ordinary gluing finishes the proof.  The argument never pulls
back an unspecified current, and never assumes that finite trace
commutes with an arbitrary base change.
\end{proof}

\begin{corollary}[The product mixed comparison in arbitrary base dimension]
\label{r15hp:product-comparison}
Let $Y=U\times\mathbf P^1$, let
$T=U\times\mathbf P^1_z\times\mathbf P^1_y$, and write
$p_z,p_y:T\to Y$ and $\pi:Y\to U$ for the indicated
projections.  For every PWS form $\gamma$ on $Y$,
\[
                 \mathcal I_{p_y}(p_z^*\gamma)
                            =\pi^*\mathcal I_\pi\gamma.
\tag{HP15.14}\label{r15hp:product-M}
\]
The left-hand representative has the prescribed current pairing
against every ordinary smooth compact complementary test on $Y$.
In particular this applies to $\gamma=\alpha$ and to
$\gamma=h\wedge\alpha$ for a smooth reference $(1,1)$ form $h$
on the projective line and any PWS $(1,1)$ input $\alpha$.
\end{corollary}

\begin{proof}
Apply Proposition~\ref{r15hp:base-change} to the algebraic
projection $S\times\mathbf P^1_y\to S$, restricted to $U$.
The all-test assertion is Theorem~\ref{r15hp:projection}, not
a test-class reduction to exterior products.
\end{proof}

\section{The fixed three Green operators over an arbitrary smooth base}
\label{r15hp:green-section}

Let $U\subset S^{\mathrm{an}}$ be as above and set
$Y=U\times\mathbf P^1$.  Retain exactly the Round 14 kernel
\[
 g=g_0+f(z)+f(y),\qquad
             Dg=h_z+h_y-\delta_\Delta,
 \qquad D=d'd'',
\tag{HP15.15}\label{r15hp:fixed-kernel}
\]
where $g_0$ is the standard chordal kernel and $f$ is its fixed
global PWS correction.  The original-field normalization of the
mass-one reference form $h$ is unchanged.

\begin{lemma}[Diagonal excision in arbitrary base dimension]
\label{r15hp:diagonal-excision}
For $\alpha\in\operatorname{PS}^{1,1}(Y)$, the three prescribed
Green current pairings can be computed locally over the base by
the PWS integrands
\[
 \chi g\,p_z^*\alpha,\qquad
 \chi d'g\wedge p_z^*\alpha,\qquad
 \chi d''g\wedge p_z^*\alpha,
\tag{HP15.16}\label{r15hp:cut-green-integrands}
\]
where $\chi$ is an ordinary PWS scalar identically zero near
$\Delta$ and one off a sufficiently small ordinary neighborhood
of $\Delta$.  The pairings are independent of all such choices
and agree on overlaps.  All compact ordinary smooth mixed tests
of the required bidegrees are included.
\end{lemma}

\begin{proof}
Along the algebraic diagonal, $p_z^*\alpha-p_y^*\alpha$
restricts to zero.  The frozen immersion-germ isomorphism makes
it zero on an ordinary neighborhood $N$ of the diagonal.
For a compact base affinoid $B\subset U$, the diagonal over
$B$ is compact and the ambient $B\times\mathbf P^1\times
\mathbf P^1$ is compact.  The chordal distance to the diagonal
therefore supplies a uniform small tube inside $N$.  A smooth
cutoff of the logarithmic chordal distance is zero on a smaller
tube and one outside this tube.  It is PWS off the diagonal
and extends as the constant zero on an ordinary neighborhood
of the diagonal.  Shrinking to an ordinary base open inside
$B$ gives the required local $\chi$.

Let $b=d+1=\dim Y$.  On the part where a cutoff changes the
integrand, replace $p_z^*\alpha$ by $p_y^*\alpha$.  For the
$G$ pairing the target test has degree $(b,b)$, so
$\alpha\wedge\eta$ has degree $(b+1,b+1)$ on $Y$ and is zero.
For $P$ the test has degree $(b-1,b)$, and the same product has
degree $(b,b+1)$ and is zero.  For $Q$ the degrees are
$(b,b-1)$ and $(b+1,b)$, with the same conclusion.  Thus even
if $d'g$ or $d''g$ has a normal component, it is wedged with
the pullback of an already zero actual form.  The paired
integrand vanishes before integration on that punctured tube.

The cutoff products in \eqref{r15hp:cut-green-integrands} extend
by zero as ordinary PWS forms across the diagonal, of degrees
$(1,1),(2,1),(1,2)$.  The first actual derivatives of the PWS
scalar $g$ away from its singular diagonal are of delta
codimension zero.  Their products with a cutoff which is zero
on an ordinary collar have genuine zero extension.  No
derivative of $\chi$ is taken.

The preceding degree calculation proves independence of any
two permitted cutoffs on every compact test.  It is the same
cutoff interpretation as the frozen prescribed Green currents;
the discarded punctured-neighborhood pairing is literally zero,
and no diagonal-supported current is added.  It also proves
compatibility for tests supported on overlapping base opens.
All remaining paired supports are compact by properness, so
the displayed integrals are defined.
\end{proof}

\begin{theorem}[Ordinary Green outputs over every smooth algebraic base]
\label{r15hp:green-outputs}
Let $S/K$ be smooth and integral, and let $U\subset S^{\mathrm{an}}$
be any ordinary open.  For every
$\alpha\in\operatorname{PS}^{1,1}(U\times\mathbf P^1)$,
without a symmetry or closedness hypothesis, the fixed prescribed
currents
\[
 \begin{split}
 G_U\alpha&=(p_y)_*[g\,p_z^*\alpha],\\
 P_U\alpha&=(p_y)_*[d'g\wedge p_z^*\alpha],\\
 Q_U\alpha&=(p_y)_*[d''g\wedge p_z^*\alpha]
 \end{split}
\tag{HP15.17}\label{r15hp:three-currents}
\]
have ordinary actual PWS representatives of degrees
$(0,0),(1,0),(0,1)$ on $U\times\mathbf P^1$.
Their pairings agree with the prescribed currents for every
ordinary smooth compact complementary test, including all mixed
tests.  They commute with ordinary restriction and with the
pullback of the base by any algebraic morphism between smooth
integral varieties, using the same fixed kernel.
\end{theorem}

\begin{proof}
On each relatively compact ordinary base open choose the three
PWS integrands \eqref{r15hp:cut-green-integrands}.
Apply Theorem~\ref{r15hp:projection} to
\[
 p_y:(U\times\mathbf P^1_y)\times\mathbf P^1_z
                                      \longrightarrow U\times\mathbf P^1_y.
\]
Here the base has dimension $d+1$ and is an ordinary open of the
smooth algebraic variety $S\times\mathbf P^1_y$.
The input degrees $(1,1),(2,1),(1,2)$ give exactly the three
claimed output degrees.  The theorem supplies genuine ordinary
PWS representatives and all smooth compact mixed-test pairings.
Lemma~\ref{r15hp:diagonal-excision} identifies these pairings
with \eqref{r15hp:three-currents}.  Current injectivity glues them
on all ordinary overlaps and removes the choice of cutoff.

There is no omitted target-collision locus in this application.
At an arbitrary target point $(x,y)$, including one where $y$
equals a root of some earlier presentation, the projection
theorem selects a new complete-fiber presentation for the
already PWS forms in \eqref{r15hp:cut-green-integrands}.
Near $z=y$ those forms are identically zero and admit the zero
presentation.  On the remaining compact source pieces, any
moment involving $z-y$ is a unit and may be perturbed within
its exact-norm class.  The resulting artificial roots are
separated at this target point.  The new target tuple consists
of genuine units; it never requires the logarithm of a
vanishing function $y-a_i(x)$.  Its use represents the same
actual integrand, not a displaced geometric diagonal.
This is why neither a bad divisor of a fixed global splitting
tuple nor a rigid-fiber reduction is required in this theorem.

For algebraic base change, choose a permitted cutoff downstairs
near the image of the point in question and pull it back to the
upstairs triple product.  The actual equality defining the
excision neighborhood pulls back, so this is also a permitted
cutoff upstairs.  Its three PWS integrands are the genuine
pullbacks of the downstairs ones.  Proposition~\ref{r15hp:base-change}
for the projective-line projection identifies their actual
outputs.  Cutoff independence and ordinary gluing prove the
asserted naturality.  No pullback of a raw current is used.
\end{proof}

\paragraph{Exact dependencies and limits.}
The new inputs are the local separable-moment perturbation,
its whole-source finite-tree presentation, and the labelled-ray
all-mixed integration theorem.  The inherited inputs are Round 14
\texttt{arbitrary\_algebraic\_immersion\_germs.tex},
\texttt{r14l:germ-iso};
\texttt{finite\_pws\_descent.tex},
\texttt{r14f:extension}, \texttt{r14f:galois-pws}, and
\texttt{r14f:pws-trace}; the separately audited native finite
degree formula; and Mihatsch Propositions 4.7 and 4.11 and
Corollary 4.9(2), as used in the labelled-ray proof.
All conclusions remain in the ordinary PWS, hence actual,
category.  No general proper actual pushforward, ambient closed
lifting theorem, or arbitrary analytic-map base-change theorem
is asserted.  Cohomology, modification descent, higher-codimension
certificates, and strong Green existence require their separate
arguments and are not consequences asserted in this section.
